\title{Large Language Models Can Automatically Engineer Features for Few-Shot Tabular Learning}

\begin{abstract}
Large Language Models (LLMs) have demonstrated exceptional capacity to address complex and novel reasoning tasks, making them highly promising for tabular learning—a core requirement in many practical settings. This work introduces \textsf{FeatLLM}, an innovative in-context learning approach where LLMs function as feature creators, crafting input datasets optimized for tabular prediction. These synthesized features enable a basic downstream machine learning model, such as linear regression, to infer class probabilities and facilitate effective few-shot learning. At inference, the \textsf{FeatLLM} framework relies exclusively on this simple model leveraging the identified features. Unlike other LLM-based strategies, \textsf{FeatLLM} dispenses with the necessity of querying the LLM for every test instance during inference, needing only API access and sidestepping prompt size restrictions. Empirical results across varied tabular datasets from multiple domains reveal that \textsf{FeatLLM} produces superior rules, outperforming existing methods like TabLLM and STUNT by a notable margin (averaging a 10\% improvement).
\end{abstract}

\section{Introduction}
The advent of Large Language Models (LLMs) has led to substantial advances in generating coherent and relevant responses to entirely new tasks, eliminating the dependency on task-specific fine-tuning~\cite{creswell2022selection,imani2023mathprompter,taylor2022galactica}. Thanks to pre-training on massive text datasets, LLMs possess comprehensive prior knowledge that can effectively replace specialized domain expertise across numerous fields~\cite{Genomes2021,singhal2023large,wilcox2003role}, thereby streamlining a variety of automation tasks, such as dialogue interpretation~\cite{finch2023leveraging} and automatic program repair~\cite{fan2023automated}. By providing clear task definitions and a handful of illustrative examples within prompts, LLMs are capable of contextualizing tasks, focusing on salient features and relevant data instances~\cite{brown2020language,zhao2023survey}. These properties highlight the potential of LLMs to serve in the role of advanced feature engineering.

Applying the knowledge base and reasoning strength of LLMs has broadened their relevance to tabular data analysis. For example, prior domain insights—such as the correlation between age and susceptibility to certain illnesses—can substantially enhance disease prediction models. Recent methodologies capture task objectives and feature details in plain language, serialize the tabular data, and input it into LLMs~\cite{dinh2022lift,wang2023anypredict}. Following this, techniques such as in-context learning~\cite{nam2023semi} or efficient fine-tuning of parameters~\cite{dinh2022lift,hegselmann2023tabllm} are employed. The incorporation of LLM-informed knowledge has shown tangible benefits, often surpassing traditional tabular learning baselines, especially in settings with limited training examples~\cite{hegselmann2023tabllm}.

Existing LLM-based approaches for tabular data learning present several drawbacks. End-to-end prediction frameworks typically demand at least one LLM inference for each data point, resulting in significant computational overhead. Moreover, most approaches depend on fine-tuning the LLM to achieve satisfactory performance, which becomes impractical for LLMs without publicly available training frameworks or APIs for adaptation—a pressing issue as state-of-the-art LLMs now often provide only constrained access through inference APIs~\cite{anil2023palm,openai2023gpt4}. Additionally, many current solutions are incompatible with scenarios requiring long prompts~\cite{liu2023lost}; as the number of features in the tabular dataset grows and the serialized input exceeds manageable lengths, these methods either fail to scale or suffer noticeable drops in predictive accuracy. Collectively, these issues pose substantial barriers to deploying LLM-based tabular models in practical contexts.

In contrast, our methodology reframes the use of LLMs in tabular settings. Rather than tasking LLMs with making direct end-to-end predictions, we leverage them as tools for feature engineering—enabling them to synthesize useful representations by capitalizing on their embedded knowledge. Specifically, we introduce \textsf{FeatLLM}, an innovative in-context learning paradigm, to discern the “criteria” underlying the predictive mechanisms of LLMs. For example, in clinical diagnosis, the LLM may identify and articulate explicit rules that associate specific feature patterns with particular disease outcomes. These extracted criteria then inform the construction of novel features, potentially supplanting original attributes for downstream modeling tasks. By repositioning the LLM as a generator of informative features, we facilitate the use of straightforward predictive models thereafter—eliminating the necessity for resource-intensive, complex models during inference and reducing latency. This feature derivation process takes advantage of LLMs’ in-context learning capabilities, permitting feature extraction by simply augmenting prompts with exemplar demonstrations rather than training the model. Further, by incorporating techniques from ensemble learning, such as feature bagging~\cite{breiman2001random}, our approach effectively addresses issues arising from excessively large prompts.

\textsf{FeatLLM} decomposes the process of generating new features from prompts into two main sub-tasks: (1) analyzing the underlying problem to infer connections between the available features and the target variable; and (2) leveraging both insights from the first stage and few-shot examples to formulate clear decision rules for class separation. This deliberate reasoning framework enables LLMs to disregard irrelevant features during rule induction. The resulting rules are executed on the dataset—interpreted as programmatic operations—transforming each sample into a binary indicator that signifies whether the rule applies. Subsequently, a simple predictive model is trained to estimate class probabilities based on these constructed features. To further enhance resilience, the approach incorporates ensembling through repeated feature generation coupled with feature bagging. By judiciously regulating the quantities of features and data samples included per bag, the method alleviates issues related to overly large prompts. Since \textsf{FeatLLM} does not require additional training and supports efficient inference, it is well-suited to diverse real-world tasks.

Our empirical evaluation involves 13 tabular datasets under few-shot conditions, demonstrating that \textsf{FeatLLM} achieves robust, high-level performance. The results highlight the LLM’s capacity to effectively leverage both domain priors and patterns in the provided data, leading to superior feature discovery. Our method consistently outperforms recent few-shot learning competitors across a range of scenarios. The implementation is made available through an anonymized GitHub repository at~\texttt{https://github.com/Sungwon-Han/FeatLLM}.

\section{Related Work}

\subsection{Few-Shot Learning with Tabular Data} Recent developments in few-shot learning have facilitated the development of robust, generalizable representations even in scenarios where annotated data is scarce~\cite{chen2018closer}. Although the initial thrust of few-shot learning targeted visual domains, the methodology has since demonstrated strong transferability to other modalities such as textual, auditory, and, most recently, tabular datasets~\cite{majumder2022few,nam2023stunt,schick2021exploiting}. However, models trained with limited labeled instances may be prone to learning non-generalizable or coincidental correlations~\cite{bartlett2021deep}. To mitigate such pitfalls, contemporary approaches increasingly rely on additional data sources or integrate domain-tailored priors to introduce meaningful inductive biases during model development. Approaches include leveraging large volumes of unlabeled tabular data, enhanced through thoughtful augmentation strategies within a semi-supervised learning framework~\cite{ucar2021subtab,yoon2020vime}, or establishing a robust foundation by means of task-agnostic unsupervised pre-training~\cite{somepalli2021saint}. Such pre-training may involve objectives that mask or intentionally corrupt segments of the input data, optimizing the model for reconstruction tasks~\cite{arik2021tabnet,majmundar2022met}, or utilize data augmentation for contrastive self-supervised learning~\cite{bahri2022scarf,chen2020simple}. Nevertheless, the intrinsic diversity within tabular datasets often renders the identification of universally effective augmentation methods problematic, and these strategies have not consistently led to improvements in few-shot environments~\cite{nam2023stunt}.

A new direction has emerged that involves training models on an expansive assortment of tabular datasets, irrespective of their connection to a specific target application, and subsequently transferring the learned knowledge to the task of interest~\cite{zhang2023generative,levin2022transfer,zhu2023xtab}. One illustrative method, TransTab~\cite{wang2022transtab}, employs transformer-based models to discern semantic interactions between columns using column metadata alongside table entries, extracting knowledge that aids zero-shot or few-shot adaptation to previously unseen tabular tasks. Another technique, TabPFN~\cite{hollmann2022tabpfn}, constructs synthetic training data by amalgamating different statistical distributions, thereby pre-training a Bayesian neural architecture. In this paradigm, inference is accomplished by providing both the support (training) and query (test) examples for the target problem directly to the pre-trained model, foregoing further adaptation. This paradigm leverages knowledge gathered across disparate datasets to outperform traditional ensemble methods, such as tree-based models, and proves particularly advantageous under few-shot conditions.

\subsection{Language-Interfaced Tabular Learning} An alternative and increasingly popular strategy for infusing prior knowledge into tabular modeling centers on the deployment of large language models (LLMs)~\cite{anil2023palm,openai2023gpt4}. LLMs, whose capabilities stem from broad pre-training over vast and heterogeneous text corpora, have exhibited remarkable generalization to novel tasks, benefiting a spectrum of disciplines~\cite{brown2020language}. There is a growing body of work exploring how to exploit the extensive knowledge housed in LLMs for tabular learning scenarios. The prevailing methodology involves the transformation of tabular inputs into text strings, which are then incorporated as part of LLM prompts~\cite{dinh2022lift,hegselmann2023tabllm,wang2023anypredict}. These prompt constructions typically consist of tabular data representations, explanatory information about features, and detailed task instructions, enabling the model to infer task-feature associations required for downstream inference. Advances in this line of research include parameter-efficient adaptation methods—such as LoRA~\cite{hu2021lora} and IA3~\cite{liu2022few}—that enable targeted fine-tuning for tabular reasoning, as well as prompt-based in-context learning techniques where a handful of exemplar instances are provided without necessitating further parameter updates~\cite{dinh2022lift,hegselmann2023tabllm,nam2023semi,slack2023tablet}.

While previous techniques have shown considerable success in enhancing few-shot tabular learning, they typically rely on routing every inference step through the LLM. This dependency creates practical obstacles, especially for deployment in industrial environments. In response, this work diverges from the standard black-box paradigm that processes each prediction in an end-to-end manner using an LLM. Rather, it prompts the LLM to directly derive class-specific rules or conditions. \textsf{FeatLLM} translates the generated rules into programmatic form to construct novel features, thus allowing downstream predictions to be performed independently of the LLM. By leveraging in-context learning, \textsf{FeatLLM} draws on both background knowledge and insights from limited-shot samples to extract interpretable rules.

\section{Methods}
The objective of \textsf{FeatLLM} is to mine the defining "rules"—namely, the characteristic conditions tied to each label—from the few available examples provided as input, as illustrated in Figure~\ref{fig:main_model}. Once obtained, these rules are parsed by an LLM, which then facilitates the generation of binary indicators for each sample, signifying whether a given rule holds. These newly formed binary features are aggregated via a weighted dot product, where the weights are non-negative and subject to training, yielding an estimate for the probability of each class. Training the model enables learning the contribution of each feature to class prediction (see Section~\ref{Sec:training}). To enhance robustness, the procedure is iterated using bagging, and predictions from different runs are amalgamated through ensembling. The subsequent sections describe the problem setting and provide comprehensive descriptions of each methodological phase.

\vspace*{-0.12in}
\paragraph{Problem formulation.} We begin with a tabular dataset containing $N$ labeled instances, expressed as $\mathcal{D} = \{(\mathbf{x}^i, \mathbf{y}^i)\}_{i=1}^{N}$. Each sample $\mathbf{x}^i$ consists of $d$ features (making it a $d$-dimensional vector), and the dataset is accompanied by a set of feature names in natural language, $F = \{f_j\}_{j=1}^{d}$, along with their respective definitions. Within the classification context, each label $\mathbf{y}^i$ belongs to a fixed set of target classes. In $k$-shot learning scenarios, a subset of $k$ ($k < N$) labeled data points is randomly selected for model training.

\subsection{Prompt Design for Extracting Rules}
\label{Sec:prompt_design}
To facilitate the LLM in deriving rules through sound reasoning processes, we steer the model's approach to reflect expert human reasoning when dealing with tabular data. The prompt is structured into three key sections (illustrated in Fig.~\ref{fig:prompt_for_rule} and detailed in Appendix~\ref{sec:example_rule}):

\vspace*{-0.12in}
\paragraph{Basic information description.} This section delivers the critical background required for the LLM to address the problem (refer to orange highlights in Fig.~\ref{fig:prompt_for_rule}). It encompasses a clear explanation of the task, information regarding the labels, comprehensive feature descriptions, and sample demonstrations. The task overview is posed as a question (for instance: ``Does this patient's myocardial infarction complications data indicate chronic heart failure? Yes or no?"; see Appendix~\ref{sec:task_desc} for additional examples). Each feature description specifies its data type and may include an explicit definition. When a feature is categorical, a list of possible categories can be added. For demonstrations, a few training examples are presented in a serialized textual format, together with their actual labels, following the serialization procedure established by earlier works~\cite{dinh2022lift,hegselmann2023tabllm}:

\begin{align}
&\text{Serialize}(\mathbf{x}^i,\mathbf{y}^i, F) = \nonumber \\ 
&``\ f_1\ \text{is}\ \mathbf{x}^i_1.\ ... \ f_d\  \text{is}\  \mathbf{x}^i_d.\ \ \text{Answer:}\  \mathbf{y}^i\ ”
\end{align}

\vspace*{-0.12in}
\paragraph{Reasoning instruction.} Our objective is to strengthen the LLM's reasoning abilities by furnishing explicit reasoning instructions (illustrated as blue text in Fig.~\ref{fig:prompt_for_rule}). The instruction commences with an opening sentence, reminiscent of the chain-of-thought methodology~\cite{wei2022chain}. We then separate the LLM's reasoning into two distinct stages. Initially, the LLM is prompted to deduce possible causal relationships or tendencies connecting features to the task description. Importantly, this initial reasoning stage is performed independently of demonstration examples, relying on general domain knowledge or common sense, thus maximizing the utilization of the model’s pre-existing understanding of the task. In the subsequent stage, the LLM is instructed to employ both the insights from the first stage and example demonstrations to induce a set of rules for each target class. This two-step structure discourages the model from learning spurious associations arising from unrelated columns and helps channel its focus towards more informative features. To control for underfitting or the risk of an excessively lengthy prompt due to generating too few or too many rules, we enforce a predefined quota of rules (specifically, 10)\footnote{Analysis on the number of rules can be found in Section~\ref{sec:analysis}.} for each class. Furthermore, while composing the rules, the LLM is instructed to consider the feature types detailed in the basic information section so as to minimize type-related mismatches.

\vspace*{-0.12in}
\paragraph{Response instruction.} To promote easier downstream processing and interpretation of the derived rules, we include response-specific instructions that specify how the LLM should organize its outputs (depicted as yellow text in Fig.~\ref{fig:prompt_for_rule}). These prompt components describe both the output structure expected for each answer class (i.e., the required response format) and the standardization of rule formulation (i.e., the rule format). By leveraging these instructions, the LLM can adapt its response template according to the type of features involved. In the absence of more restrictive instructions, the LLM retains the flexibility to combine individual rules using logical connectors such as AND or OR. For reference, a sample of the resultant output appears in Appendix~\ref{sec:outcome-rule}.

\subsection{Inferring Class Likelihood via Rules}
\label{Sec:training}

\paragraph{Parsing rules for feature generation.} The rules derived in the earlier phase serve as the basis for constructing new binary features. For every class, a feature is generated that reflects whether a specific sample meets that class's rules, taking the value ‘0’ or ‘1’. These features assist in estimating the probability that a sample pertains to each class. Because the rules obtained from the LLM are written in natural language, an effective parsing approach is essential for automated feature construction. Although we impose formatting constraints and specify rules clearly to ease parsing during rule generation, the LLM’s responses often exhibit syntactic noise~\cite{kovavc2023large}, such as rephrasing operators (e.g., using “greater than” or “less than” instead of the symbols “$>$” or “$<$”). Such inconsistencies complicate parsing.

To cope with noisy, syntactically varied outputs, we forego intricate programmatic parsing and instead employ the LLM itself for parsing. The LLM excels at semantic comprehension across syntactic variations and can generate functional code with minimal guidance, owing to its exposure to code in its training. To leverage this, we craft prompts containing the function name, detailed descriptions of input and output, and the identified rules, all of which are supplied to the LLM. Examples of these prompts and their corresponding outputs are provided in Figures~\ref{fig:prompt_for_parsing} and~\ref{sec:example_parsing}-\ref{sec:example_parse_outcome} in the Appendix. The code generated from the LLM is subsequently executed via Python’s \texttt{exec()} to carry out the data transformation.

\vspace*{-0.12in}
\paragraph{Inferring class likelihood.} With these newly derived binary features per class, one straightforward approach for estimating the likelihood that a sample belongs to a class is to compute the number of class-specific rules it satisfies (that is, the aggregate of the binary features for a given class). Nonetheless, the relevance of individual rules and features can differ, which necessitates learning their respective contributions from labeled data. We propose to estimate these importances by training a standard machine learning (ML) model using the binary features of each class. For instance, a linear model without a bias term can be utilized. Specifically, let the binary feature vector for sample $\mathbf{x}^i$ with respect to class $k$ be denoted as $\mathbf{z}_k^i \in \{0, 1\}^R$, where $R$ stands for the rule count per class and $c$ denotes the number of classes. The probability assigned to each class, $\mathbf{p}^i$, is then determined as follows:

\begin{align}
\text{logit}_k^i &= \max(\mathbf{w}_k, 0) \cdot \mathbf{z}_k^i, \nonumber \\
\mathbf{p}^i &= \text{Softmax}({[\text{logit}_{1}^i, ..., \text{logit}_{c}^i]}). \label{eq:infer_prob}
\end{align}

In this context, $\mathbf{w}_k$ denotes the adjustable parameters in the linear classifier, capturing the relevance of individual features. By constraining these weights to the non-negative domain, we guarantee that the outputs generated by the LLM are prevented from negatively affecting the class probabilities during model training. The resulting logits are subsequently mapped to probability distributions over classes via the softmax transformation. We fit $\mathbf{w}_k$ by minimizing cross-entropy loss, utilizing a small set of labeled examples in a few-shot regime.

\vspace*{-0.12in}
\paragraph{Ensembling with bagging.} In the final stage, we perform the entire inference process multiple times, generating an ensemble of prediction models whose outputs are combined for the final decision. Specifically, after obtaining trained weight vectors from $T$ separate runs, denoted by $\{\mathbf{w}[t]\}_{t=1}^T$, we calculate each class probability $\mathbf{p}^i$ for all $t$ using Eq.~\ref{eq:infer_prob} and aggregate them by averaging.

To promote variability in the ensemble’s rule sets, we configure the LLM with a high temperature during sampling. We also shuffle the order of the provided few-shot demonstration examples, since prior work has demonstrated that this order can influence LLM outputs~\cite{lu2022fantastically}\footnote{Further discussion of rule diversity appears in Appendix~\ref{sec:diversity}}. In addition, bagging is applied by randomly sampling subsets of features or data instances for each ensemble member, a practice especially advantageous when handling larger training sets or feature spaces. This ensemble strategy brings two principal benefits. First, the effect of erroneously induced rules in some trials is mitigated by the presence of accurate rules in others, thereby improving the system’s resilience; this draws on the LLM’s self-consistency property~\cite{wang2022self}, since repeatedly generated rules tend to be more reliable. Second, ensembling helps address prompt length constraints inherent to LLMs. When tabular datasets surpass prompt capacity, bagging can restrict the input size, maintaining dependable predictive performance. Thus, while an individual set of rules might miss some feature relations, assembling diverse weak models across trials enables the collective ensemble to overcome the limitations of incomplete feature or sample coverage within single runs.

\section{Experiments}
We assess the effectiveness of \textsf{FeatLLM} using a variety of tabular datasets, and carry out multiple analyses to understand the inner workings of our framework. Further experimental results, including tests with different LLMs, qualitative examinations, and more, can be found in the Appendix due to space limitations.

\subsection{Performance Evaluation}
\paragraph{Datasets.}
Our experimental setup draws on 13 datasets encompassing binary and multi-class classification problems: (1) Adult~\cite{asuncion2007uci}, targeting the prediction of whether a person’s annual income surpasses \$50,000; (2) Bank~\cite{moro2014data}, for forecasting whether a customer will opt in for a term deposit; (3) Blood~\cite{yeh2009knowledge}, predicting the likelihood of donors returning for future donations; (4) Car~\cite{kadra2021well}, focused on assessing the quality level of vehicles; (5) Communities~\cite{misc_communities_and_crime_183}, for estimating crime rates within particular regions; (6) Credit-g~\cite{kadra2021well}, predicting if an individual presents a good or bad credit risk; (7) Diabetes\footnote{\texttt{kaggle.com/datasets/uciml/pima-indians-diabetes-database}}, determining the presence of diabetes in a given subject; (8) Heart\footnote{\texttt{kaggle.com/datasets/fedesoriano/heart-failure-prediction}}, forecasting the development of coronary artery disease; and (9) Myocardial~\cite{misc_myocardial_infarction_complications_579}, identifying individuals suffering from chronic heart failure.

Moreover, we integrate newer and synthetic datasets that are underexplored in literature and are absent from the pre-training corpus of LLMs: (10) Cultivars, which aims to estimate the grain yield level for a soybean cultivar\footnote{\texttt{archive.ics.uci.edu/dataset/913}}; (11) NHANES, predicting an individual’s age bracket from their record\footnote{\texttt{archive.ics.uci.edu/dataset/887}}; (12) Sequence-type, a classification task distinguishing between four sequence types: arithmetic, geometric, Fibonacci, or Collatz; (13) Solution-mix, where the goal is to identify if the concentration percentage in a mixture from four solutions exceeds 0.5. The first two were only recently made available in the UCI Machine Learning Repository post-September 2023, which guarantees that they do not overlap with the pre-training set of the employed LLM API (gpt-3.5-turbo-0613). The latter two are synthetic, precluding any chance of memorization by the LLM and supporting a fair evaluation.

A summary of dataset properties such as scale and intricacy is provided in Table~\ref{Tab:basic-desc}.
Every dataset includes well-defined attribute names and descriptions. Particularly, datasets such as `Communities' and `Myocardial' encompass over 100 features, making them especially challenging because of the context window limitations in LLMs.

\vspace*{-0.12in}
\paragraph{Baselines.}
We benchmark our framework against nine existing approaches. The first group comprises standard tabular modeling techniques: (1) Logistic regression (LogReg), (2) XGBoost~\cite{chen2016xgboost}, and (3) RandomForest~\cite{ho1995random}. The next two—(4) SCARF~\cite{bahri2022scarf} and (5) STUNT~\cite{nam2023stunt}—assume access to large volumes of unlabeled data, an assumption that may not hold in practice. Another relevant baseline is (6) TabPFN~\cite{hollmann2022tabpfn}, which relies on synthetic, diverse data distributions for pretraining. The last three approaches utilize LLMs, namely: (7) In-context learning (In-context)~\cite{wei2022emergent}, which inserts few-shot examples directly into the prompt without parameter adaptation; (8) TABLET~\cite{slack2023tablet}, augmenting prompts with external classifier information—such as rule-sets and prototypes—to bolster inference; and (9) TabLLM~\cite{hegselmann2023tabllm}, which adopts parameter-efficient tuning (e.g., IA3~\cite{liu2022few}) to train the LLM using limited samples.

\vspace*{-0.12in}
\paragraph{Implementation details.}
\textsf{FeatLLM} is built upon the GPT-3.5 architecture as its primary LLM, but it maintains compatibility with other LLM variants (refer to Appendix~\ref{sec:backbones} for an analysis of alternate backbones). The inference parameters are configured with a temperature of 0.5 and a top-p of 1, in line with API defaults. For rule extraction, we set the ensemble size and number of rules to 20 and 10, respectively. An exploration of hyper-parameter sensitivity can be found in Figure~\ref{fig:hyperparameter2} and detailed further in Appendix~\ref{sec:hyper-parameter}. 

Training the linear component employs the Adam optimizer with a learning rate of 0.01 for 200 epochs, selecting the optimal epoch via $k$-fold cross-validation. For the baseline models, in-context learning is conducted with GPT-3.5, while TabLLM utilizes T0~\cite{sanh2022multitask} to match the original configuration. Importantly, TabLLM is limited to LLMs that have publicly available checkpoints, excluding GPT-3.5. As for traditional models such as LogReg, XGBoost, and RandomForest, parameter optimization is achieved through a combination of Grid-search and $k$-fold cross-validation, where $k$ is chosen as 2 or 4 to ensure that every class is present in the training split. Further implementation specifics can be found in Appendix~\ref{sec:baselines}.

\vspace*{-0.12in}
\paragraph{Results.}
Tables~\ref{tab:main1} and \ref{tab:main2} report comparative performance results over various datasets. Each experiment was conducted three times, with the selection of training and test splits randomized, and we report both the mean and standard deviation of the resulting AUC metrics. Across all benchmarks, \textsf{FeatLLM} is either the highest performer or ranks second, outperforming most baseline approaches. Remarkably, with only 4 shots, our method matches the efficacy of standard tabular baselines that need 32 shots (refer to Fig.~\ref{fig:compares}). Although models like STUNT and TabPFN achieved notable gains on individual datasets, their aggregate improvements over classic machine learning algorithms remained modest. Additionally, LLM-based methods—including in-context learning, TABLET, and TabLLM—were unable to process tabular datasets featuring a large quantity of variables. By incorporating bagging and ensemble methods, our framework effectively handles datasets with many features, maintaining competitive results. Notably, this strategy of bagging and ensembling could, in principle, be extended to other LLM-centered models; however, doing so would necessitate twice as many per-sample inferences, significantly elevating computational demands and diminishing practical applicability.

\subsection{Analysis \& Discussion}
\label{sec:analysis}
\paragraph{Ablation study.} To assess the individual impact of primary system components, we conducted ablation experiments evaluating: (1) \textsf{-Tuning}, where the linear model's weight tuning step is removed and binary features are simply summed to generate logits; (2) \textsf{-Ensemble}, where ensemble aggregation is omitted; (3) \textsf{-Description}, which leaves out feature descriptions; and (4) \textsf{-Reasoning}, which excludes the Step-1 instruction during rule generation.

The results, summarized in Table~\ref{tab:ablation}, present the average changes in AUC across all datasets when each component is ablated. In every scenario, omitting a given part leads to a decrease in performance. Notably, the utility of weight tuning escalates with an increasing number of shots, indicating that greater data availability enhances the reliable assessment of rule importance, thereby making tuning particularly impactful. In contrast, the value of feature description and reasoning instructions is most pronounced at lower shot counts, implying that leveraging LLM’s prior knowledge is especially important in data-scarce situations. Finally, the ensemble method introduced consistently yields improvements in every experimental configuration.

\vspace*{-0.12in}
\paragraph{Training \& inference time comparison.} We assess both the training and inference durations across various methods, conducting all experiments using the Adult dataset with a 16-shot training split. All models were executed on a single A100 GPU by default; the exception is TabLLM, which necessitates four A100 GPUs for model parallelization. Table~\ref{tab:cost} presents a detailed comparison of runtimes. Remarkably, \textsf{FeatLLM} delivers inference times that are quite close to those of traditional approaches. In stark contrast, inference times for alternative LLM-based techniques—including In-context learning, TABLET, and TabLLM—are substantially higher. Regarding training overhead, \textsf{FeatLLM} remains on par with other LLM-oriented approaches, requiring merely 30 API calls, which minimally impacts resource budgets and is amenable to parallel processing.

\vspace*{-0.12in}
\paragraph{Quality of features generated by \textsf{FeatLLM}.}
To evaluate the informativeness of the rule-based features generated by \textsf{FeatLLM}, we perform a straightforward empirical test. For each dataset, we calculate the mean AUROC between the generated features and the target labels, then determine the proportion of features deemed useful—those yielding an AUROC above 0.5. Table~\ref{tab:quality} details these findings. The results demonstrate that a substantial portion of the constructed rules bear significant relevance to the target label and enrich the feature set for downstream tasks (see the ``Before tuning'' column). Further, when the linear model is subsequently fine-tuned on data, irrelevant or low-utility rules (those whose learned coefficients fall beneath a set threshold) are systematically eliminated (as shown in the ``After tuning'' column). For these experiments, a threshold value of 0.1 was adopted, based on the distribution of learned weights.

\vspace*{-0.12in}
\paragraph{Effect of adding spurious correlations.} To evaluate the capability of different approaches in eliminating irrelevant spurious correlations under data-scarce conditions, we performed a targeted experiment. Specifically, we selected the Heart dataset and augmented it by appending randomly chosen columns from the Adult dataset, ensuring that these columns bore no relevance to the Heart task. The model performance was then tracked as the number of these unrelated columns increased. For consistency across comparisons, each method was provided with 8 labeled samples only, and no unlabeled data was utilized, thus omitting approaches such as SCARF and STUNT from this evaluation.

Figure~\ref{fig:spurious} presents how model performance shifts when spurious correlations are introduced. Among all tested methods, \textsf{FeatLLM} exhibited the smallest reduction in performance as additional irrelevant columns were added. This resilience can be attributed to our framework’s strategy, which compels the extraction of rules anchored in domain-specific prior knowledge about feature-task associations, thereby minimizing the risk of incorporating rules based on irrelevant columns. As a result, rules referencing the noisy columns represented merely 1-2\% of the overall rule set. Nonetheless, it is important to note that \textsf{FeatLLM} is not immune to learning from spurious correlations, particularly when extraneous columns come from datasets sharing conceptual similarities with the target task (further evidence is discussed in Appendix~\ref{sec:spurious-appendix}).

\vspace*{-0.12in}
\paragraph{Hyper-parameter analysis}
\textsf{FeatLLM} incorporates two principal hyper-parameters: the number of ensemble models utilized, and the count of rules extracted per class during each interaction with the LLM. To assess the influence of these hyper-parameters on system performance, we conducted a series of experiments. The findings indicate that increasing the ensemble size generally leads to enhanced performance; however, this improvement plateaus after approximately 20 ensembles (refer to Fig.~\ref{fig:appendix-hyperparameter1} in the Appendix). Regarding the number of rules, while no statistically significant variations were observed, selecting 10 rules was found to offer the optimal compromise between prompt complexity and model accuracy (see Fig.~\ref{fig:hyperparameter2}). We hypothesize that although augmenting the number of rules can enrich the information content by expanding input dimensionality, it may also exacerbate overfitting, especially in scenarios with limited data.

\section{Conclusion}
We introduce \textsf{FeatLLM}, which establishes a new benchmark in LLM-driven few-shot tabular data learning. Unlike conventional methods that rely on LLMs solely for sample-level predictions, \textsf{FeatLLM} leverages these models for generating explicit decision criteria, akin to feature engineering. By combining rule extraction with an ensemble approach utilizing bagging, \textsf{FeatLLM} achieves high predictive accuracy with minimal inference requirements, operates through API access without necessitating model training, and circumvents restrictions on data dimensionality. As a result, \textsf{FeatLLM} emerges as an efficient and practical solution for real-world tabular datasets, delivering strong performance in a data-conservative manner.

Currently, \textsf{FeatLLM} targets tasks characterized by scarce labeled data. Moving forward, we aim to broaden its applicability to larger datasets by enabling the generation of more diverse feature types beyond simple rules, enhancing interpretability, and thereby extending its utility across a broader spectrum of applications.

\section{Broader Impact}
\textsf{FeatLLM} is conceived to capitalize on LLMs' extensive prior knowledge, pushing beyond the standard limitations encountered in traditional supervised tabular learning. By augmenting learning processes with such prior knowledge, our work substantially contributes to reducing overfitting when working with small sample sizes, thereby improving data efficiency. This approach holds significant value for professionals in industries like finance and healthcare, where labeling data can be prohibitively costly or complex, and LLMs—already pre-trained on domain-relevant corpora—can provide crucial auxiliary information. Nevertheless, for widespread deployment, it is critical to systematically examine how societal biases or inaccuracies inherent in LLM pre-training may impact downstream predictions, as these may heavily influence or even overshadow empirical evidence from collected labels.

\end{document}